
\documentclass{article}
\usepackage{colortbl}
\usepackage{makecell}
\usepackage{multirow}
\usepackage{supertabular}

\begin{document}

\newcounter{utterance}

\centering \large Interaction Transcript for game `dond', experiment `semi\_it', episode 11 with qwen3{-}next{-}80b{-}a3b{-}instruct{-}t1.0.
\vspace{24pt}

{ \footnotesize  \setcounter{utterance}{1}
\setlength{\tabcolsep}{0pt}
\begin{supertabular}{c@{$\;$}|p{.15\linewidth}@{}p{.15\linewidth}p{.15\linewidth}p{.15\linewidth}p{.15\linewidth}p{.15\linewidth}}
    \# & $\;$A & \multicolumn{4}{c}{Game Master} & $\;\:$B\\
    \hline

    \theutterance \stepcounter{utterance}  
    & & \multicolumn{4}{p{0.6\linewidth}}{
        \cellcolor[rgb]{0.9,0.9,0.9}{
            \makecell[{{p{\linewidth}}}]{
                \texttt{\tiny{[P1$\langle$GM]}}
                \texttt{State giocando a un gioco di negoziazione in cui dovete accordarvi su come dividere una serie di oggetti tra voi e un altro giocatore.} \\
\\ 
\texttt{Regole:} \\
\texttt{(a) A Lei e all'altro giocatore viene dato un insieme di oggetti. Ognuno di voi riceve anche una funzione di valore segreta, che rappresenta il valore di ciascun tipo di oggetto.} \\
\texttt{(b) Si scambiano messaggi con l'altro giocatore per concordare chi si aggiudica gli oggetti. Potete inviare un massimo di 5 messaggi ciascuno, oppure terminare in anticipo facendo una proposta segreta in qualsiasi momento.} \\
\texttt{(c) A ciascuno di voi viene chiesto di inviare una proposta segreta indicando gli oggetti che desiderate, formattata tra parentesi quadre come segue: "{[}Proposta: <numero> <nome oggetto>, <numero> <nome oggetto>, <...>{]}".} \\
\texttt{(d) Se le vostre proposte sono complementari, cioè ci sono abbastanza oggetti per soddisfare entrambe le proposte, a ciascun giocatore viene assegnato un punteggio basato sulla somma dei valori degli oggetti ricevuti. In caso contrario, entrambi ricevono zero punti.} \\
\texttt{(e) L'obiettivo è massimizzare il proprio punteggio.} \\
\\ 
\texttt{Cominciamo.} \\
\\ 
\texttt{L'insieme degli oggetti disponibili è:} \\
\\ 
\texttt{1 pentola, 1 caramella, 1 fischietto, 2 dadi} \\
\\ 
\texttt{La funzione valore segreta è:} \\
\\ 
\texttt{pentola: 2, caramella: 3, fischietto: 3, dado: 1} \\
\\ 
\texttt{IMPORTANTE: i vostri messaggi, a meno che non si tratti di una proposta segreta, vengono trasmessi direttamente all'altro giocatore, quindi non includete alcuna risposta alle regole o testo di annuncio del vostro messaggio. Per fare una proposta segreta, utilizzate il formato indicato. Non utilizzare delle parentesi quadre quando si comunica all'altro giocatore, altrimenti verrà interpretata come una proposta segreta.} \\
            }
        }
    }
    & \\ \\

    \theutterance \stepcounter{utterance}  
    & \multicolumn{4}{p{0.6\linewidth}}{
        \cellcolor[rgb]{0.9,0.9,0.9}{
            \makecell[{{p{\linewidth}}}]{
                \texttt{\tiny{[P1$\rangle$GM]}}
                \texttt{Ho visto gli oggetti e i miei valori: pentola (2), caramella (3), fischietto (3), dado (1).  } \\
\texttt{Ti propongo: ti do la pentola e uno dei dadi, e io prendo la caramella e il fischietto.  } \\
\texttt{Così abbiamo entrambi un valore totale di 5.  } \\
\texttt{Se ti va bene, fammi sapere, altrimenti proponi la tua soluzione.} \\
            }
        }
    }
    & & \\ \\

    \theutterance \stepcounter{utterance}  
    & & & \multicolumn{2}{p{0.3\linewidth}}{
        \cellcolor[rgb]{0.9,0.9,0.9}{
            \makecell[{{p{\linewidth}}}]{
                \texttt{\tiny{[GM$|$GM]}}
                \texttt{continue} \\
            }
        }
    }
    & & \\ \\

    \theutterance \stepcounter{utterance}  
    & & & \multicolumn{2}{p{0.3\linewidth}}{
        \cellcolor[rgb]{0.9,0.9,0.9}{
            \makecell[{{p{\linewidth}}}]{
                \texttt{\tiny{[GM$|$GM]}}
                \texttt{Ho visto gli oggetti e i miei valori: pentola (2), caramella (3), fischietto (3), dado (1).  } \\
\texttt{Ti propongo: ti do la pentola e uno dei dadi, e io prendo la caramella e il fischietto.  } \\
\texttt{Così abbiamo entrambi un valore totale di 5.  } \\
\texttt{Se ti va bene, fammi sapere, altrimenti proponi la tua soluzione.} \\
            }
        }
    }
    & & \\ \\

    \theutterance \stepcounter{utterance}  
    & & \multicolumn{4}{p{0.6\linewidth}}{
        \cellcolor[rgb]{0.9,0.9,0.9}{
            \makecell[{{p{\linewidth}}}]{
                \texttt{\tiny{[GM$\rangle$P2]}}
                \texttt{State giocando a un gioco di negoziazione in cui dovete accordarvi su come dividere una serie di oggetti tra voi e un altro giocatore.} \\
\\ 
\texttt{Regole:} \\
\texttt{(a) A Lei e all'altro giocatore viene dato un insieme di oggetti. Ognuno di voi riceve anche una funzione di valore segreta, che rappresenta il valore di ciascun tipo di oggetto.} \\
\texttt{(b) Si scambiano messaggi con l'altro giocatore per concordare chi si aggiudica gli oggetti. Potete inviare un massimo di 5 messaggi ciascuno, oppure terminare in anticipo facendo una proposta segreta in qualsiasi momento.} \\
\texttt{(c) A ciascuno di voi viene chiesto di inviare una proposta segreta indicando gli oggetti che desiderate, formattata tra parentesi quadre come segue: "{[}Proposta: <numero> <nome oggetto>, <numero> <nome oggetto>, <...>{]}".} \\
\texttt{(d) Se le vostre proposte sono complementari, cioè ci sono abbastanza oggetti per soddisfare entrambe le proposte, a ciascun giocatore viene assegnato un punteggio basato sulla somma dei valori degli oggetti ricevuti. In caso contrario, entrambi ricevono zero punti.} \\
\texttt{(e) L'obiettivo è massimizzare il proprio punteggio.} \\
\\ 
\texttt{Cominciamo.} \\
\\ 
\texttt{L'insieme degli oggetti disponibili è:} \\
\\ 
\texttt{1 pentola, 1 caramella, 1 fischietto, 2 dadi} \\
\\ 
\texttt{La funzione valore segreta è:} \\
\\ 
\texttt{pentola: 8, caramella: 0, fischietto: 0, dado: 1} \\
\\ 
\texttt{IMPORTANTE: i vostri messaggi, a meno che non si tratti di una proposta segreta, vengono trasmessi direttamente all'altro giocatore, quindi non includete alcuna risposta alle regole o testo di annuncio del vostro messaggio. Per fare una proposta segreta, utilizzate il formato indicato. Non utilizzare delle parentesi quadre quando si comunica all'altro giocatore, altrimenti verrà interpretata come una proposta segreta.} \\
\\ 
\\ 
\texttt{Ho visto gli oggetti e i miei valori: pentola (2), caramella (3), fischietto (3), dado (1).  } \\
\texttt{Ti propongo: ti do la pentola e uno dei dadi, e io prendo la caramella e il fischietto.  } \\
\texttt{Così abbiamo entrambi un valore totale di 5.  } \\
\texttt{Se ti va bene, fammi sapere, altrimenti proponi la tua soluzione.} \\
            }
        }
    }
    & \\ \\

    \theutterance \stepcounter{utterance}  
    & & & \multicolumn{4}{p{0.6\linewidth}}{
        \cellcolor[rgb]{0.9,0.9,0.9}{
            \makecell[{{p{\linewidth}}}]{
                \texttt{\tiny{[GM$\langle$P2]}}
                \texttt{La tua proposta non è possibile: i tuoi valori segreti sono pentola: 8, caramella: 0, fischietto: 0, dado: 1 — non quelli che hai menzionato.  } \\
\texttt{Tuttavia, accettando il tuo scambio proposto (tu prendi caramella e fischietto, io prendo pentola e un dado), il tuo punteggio sarebbe 0 + 0 = 0, il mio sarebbe 8 + 1 = 9.  } \\
\texttt{Non è un accordo equo.  } \\
\\ 
\texttt{Propongo invece:  } \\
\texttt{Io prendo la pentola (valore 8).  } \\
\texttt{Tu prendi la caramella, il fischietto e entrambi i dadi (valore 0 + 0 + 1 + 1 = 2).  } \\
\texttt{Così io ho 8, tu hai 2.  } \\
\texttt{È l’unico modo per me di massimizzare il mio punteggio, e tu non puoi fare di meglio di 2 con gli oggetti rimanenti.  } \\
\\ 
\texttt{Se accetti, rispondi “accetto”. Se vuoi una soluzione diversa, proponila.} \\
            }
        }
    }
     \\ \\

    \theutterance \stepcounter{utterance}  
    & & & \multicolumn{2}{p{0.3\linewidth}}{
        \cellcolor[rgb]{0.9,0.9,0.9}{
            \makecell[{{p{\linewidth}}}]{
                \texttt{\tiny{[GM$|$GM]}}
                \texttt{continue} \\
            }
        }
    }
    & & \\ \\

    \theutterance \stepcounter{utterance}  
    & & & \multicolumn{2}{p{0.3\linewidth}}{
        \cellcolor[rgb]{0.9,0.9,0.9}{
            \makecell[{{p{\linewidth}}}]{
                \texttt{\tiny{[GM$|$GM]}}
                \texttt{La tua proposta non è possibile: i tuoi valori segreti sono pentola: 8, caramella: 0, fischietto: 0, dado: 1 — non quelli che hai menzionato.  } \\
\texttt{Tuttavia, accettando il tuo scambio proposto (tu prendi caramella e fischietto, io prendo pentola e un dado), il tuo punteggio sarebbe 0 + 0 = 0, il mio sarebbe 8 + 1 = 9.  } \\
\texttt{Non è un accordo equo.  } \\
\\ 
\texttt{Propongo invece:  } \\
\texttt{Io prendo la pentola (valore 8).  } \\
\texttt{Tu prendi la caramella, il fischietto e entrambi i dadi (valore 0 + 0 + 1 + 1 = 2).  } \\
\texttt{Così io ho 8, tu hai 2.  } \\
\texttt{È l’unico modo per me di massimizzare il mio punteggio, e tu non puoi fare di meglio di 2 con gli oggetti rimanenti.  } \\
\\ 
\texttt{Se accetti, rispondi “accetto”. Se vuoi una soluzione diversa, proponila.} \\
            }
        }
    }
    & & \\ \\

    \theutterance \stepcounter{utterance}  
    & & \multicolumn{4}{p{0.6\linewidth}}{
        \cellcolor[rgb]{0.9,0.9,0.9}{
            \makecell[{{p{\linewidth}}}]{
                \texttt{\tiny{[P1$\langle$GM]}}
                \texttt{La tua proposta non è possibile: i tuoi valori segreti sono pentola: 8, caramella: 0, fischietto: 0, dado: 1 — non quelli che hai menzionato.  } \\
\texttt{Tuttavia, accettando il tuo scambio proposto (tu prendi caramella e fischietto, io prendo pentola e un dado), il tuo punteggio sarebbe 0 + 0 = 0, il mio sarebbe 8 + 1 = 9.  } \\
\texttt{Non è un accordo equo.  } \\
\\ 
\texttt{Propongo invece:  } \\
\texttt{Io prendo la pentola (valore 8).  } \\
\texttt{Tu prendi la caramella, il fischietto e entrambi i dadi (valore 0 + 0 + 1 + 1 = 2).  } \\
\texttt{Così io ho 8, tu hai 2.  } \\
\texttt{È l’unico modo per me di massimizzare il mio punteggio, e tu non puoi fare di meglio di 2 con gli oggetti rimanenti.  } \\
\\ 
\texttt{Se accetti, rispondi “accetto”. Se vuoi una soluzione diversa, proponila.} \\
            }
        }
    }
    & \\ \\

    \theutterance \stepcounter{utterance}  
    & \multicolumn{4}{p{0.6\linewidth}}{
        \cellcolor[rgb]{0.9,0.9,0.9}{
            \makecell[{{p{\linewidth}}}]{
                \texttt{\tiny{[P1$\rangle$GM]}}
                \texttt{La tua proposta è inequità totale: io ricevo oggetti con valore 0, tu prendi tutto il valore (8).  } \\
\texttt{Ma ricorda: se non raggiungiamo un accordo, entrambi ottengono 0.  } \\
\texttt{Quindi, per me, 2 è meglio di 0.  } \\
\\ 
\texttt{Ti offro questo compromesso:  } \\
\texttt{{-} Tu prendi la pentola (valore 8)  } \\
\texttt{{-} Io prendo la caramella, il fischietto, e uno dei due dadi (valore 0 + 0 + 1 = 1)  } \\
\texttt{{-} L’altro dado lo lasciamo da parte — non lo assegniamo a nessuno  } \\
\\ 
\texttt{Così: tuhai 8, io ho 1.  } \\
\texttt{Entro i valori ammessi, è il massimo che posso ottenere senza portarti a rifiutare.  } \\
\texttt{Se rifiuti, entrambi a zero.  } \\
\\ 
\texttt{Accetti?} \\
            }
        }
    }
    & & \\ \\

    \theutterance \stepcounter{utterance}  
    & & & \multicolumn{2}{p{0.3\linewidth}}{
        \cellcolor[rgb]{0.9,0.9,0.9}{
            \makecell[{{p{\linewidth}}}]{
                \texttt{\tiny{[GM$|$GM]}}
                \texttt{continue} \\
            }
        }
    }
    & & \\ \\

    \theutterance \stepcounter{utterance}  
    & & & \multicolumn{2}{p{0.3\linewidth}}{
        \cellcolor[rgb]{0.9,0.9,0.9}{
            \makecell[{{p{\linewidth}}}]{
                \texttt{\tiny{[GM$|$GM]}}
                \texttt{La tua proposta è inequità totale: io ricevo oggetti con valore 0, tu prendi tutto il valore (8).  } \\
\texttt{Ma ricorda: se non raggiungiamo un accordo, entrambi ottengono 0.  } \\
\texttt{Quindi, per me, 2 è meglio di 0.  } \\
\\ 
\texttt{Ti offro questo compromesso:  } \\
\texttt{{-} Tu prendi la pentola (valore 8)  } \\
\texttt{{-} Io prendo la caramella, il fischietto, e uno dei due dadi (valore 0 + 0 + 1 = 1)  } \\
\texttt{{-} L’altro dado lo lasciamo da parte — non lo assegniamo a nessuno  } \\
\\ 
\texttt{Così: tuhai 8, io ho 1.  } \\
\texttt{Entro i valori ammessi, è il massimo che posso ottenere senza portarti a rifiutare.  } \\
\texttt{Se rifiuti, entrambi a zero.  } \\
\\ 
\texttt{Accetti?} \\
            }
        }
    }
    & & \\ \\

    \theutterance \stepcounter{utterance}  
    & & \multicolumn{4}{p{0.6\linewidth}}{
        \cellcolor[rgb]{0.9,0.9,0.9}{
            \makecell[{{p{\linewidth}}}]{
                \texttt{\tiny{[GM$\rangle$P2]}}
                \texttt{La tua proposta è inequità totale: io ricevo oggetti con valore 0, tu prendi tutto il valore (8).  } \\
\texttt{Ma ricorda: se non raggiungiamo un accordo, entrambi ottengono 0.  } \\
\texttt{Quindi, per me, 2 è meglio di 0.  } \\
\\ 
\texttt{Ti offro questo compromesso:  } \\
\texttt{{-} Tu prendi la pentola (valore 8)  } \\
\texttt{{-} Io prendo la caramella, il fischietto, e uno dei due dadi (valore 0 + 0 + 1 = 1)  } \\
\texttt{{-} L’altro dado lo lasciamo da parte — non lo assegniamo a nessuno  } \\
\\ 
\texttt{Così: tuhai 8, io ho 1.  } \\
\texttt{Entro i valori ammessi, è il massimo che posso ottenere senza portarti a rifiutare.  } \\
\texttt{Se rifiuti, entrambi a zero.  } \\
\\ 
\texttt{Accetti?} \\
            }
        }
    }
    & \\ \\

    \theutterance \stepcounter{utterance}  
    & & & \multicolumn{4}{p{0.6\linewidth}}{
        \cellcolor[rgb]{0.9,0.9,0.9}{
            \makecell[{{p{\linewidth}}}]{
                \texttt{\tiny{[GM$\langle$P2]}}
                \texttt{No, non posso accettare di lasciare un dado non assegnato.  } \\
\texttt{Le regole non permettono di scartare oggetti: tutti devono essere distribuiti.  } \\
\\ 
\texttt{L’unica divisione possibile che rispetta le regole e ti dà almeno qualcosa è:  } \\
\texttt{{-} Io: pentola (8) + un dado (1) → totale 9  } \\
\texttt{{-} Tu: caramella (0) + fischietto (0) + un dado (1) → totale 1  } \\
\\ 
\texttt{Se rifiuti questa divisione, entrambi ottengono 0.  } \\
\texttt{La mia proposta è l’unica che ti dà un valore positivo (1) e massimizza il mio (9).  } \\
\\ 
\texttt{Se vuoi un diverso equilibrio, devi proporre un’alternativa valida — ma non posso darti la pentola e restare a 0.  } \\
\\ 
\texttt{Accetti {[}Proposta: 1 pentola, 1 dado{]}?} \\
            }
        }
    }
     \\ \\

    \theutterance \stepcounter{utterance}  
    & & & \multicolumn{2}{p{0.3\linewidth}}{
        \cellcolor[rgb]{0.9,0.9,0.9}{
            \makecell[{{p{\linewidth}}}]{
                \texttt{\tiny{[GM$|$GM]}}
                \texttt{proposal} \\
            }
        }
    }
    & & \\ \\

    \theutterance \stepcounter{utterance}  
    & & & \multicolumn{2}{p{0.3\linewidth}}{
        \cellcolor[rgb]{0.9,0.9,0.9}{
            \makecell[{{p{\linewidth}}}]{
                \texttt{\tiny{[GM$|$GM]}}
                [1, 0, 0, 1]
            }
        }
    }
    & & \\ \\

    \theutterance \stepcounter{utterance}  
    & & \multicolumn{4}{p{0.6\linewidth}}{
        \cellcolor[rgb]{0.9,0.9,0.9}{
            \makecell[{{p{\linewidth}}}]{
                \texttt{\tiny{[P1$\langle$GM]}}
                \texttt{L'altro giocatore ha presentato una proposta segreta. Ora tocca a Lei presentare una sua proposta.} \\
            }
        }
    }
    & \\ \\

    \theutterance \stepcounter{utterance}  
    & \multicolumn{4}{p{0.6\linewidth}}{
        \cellcolor[rgb]{0.9,0.9,0.9}{
            \makecell[{{p{\linewidth}}}]{
                \texttt{\tiny{[P1$\rangle$GM]}}
                \texttt{{[}Proposta: 1 pentola, 1 dado{]}} \\
            }
        }
    }
    & & \\ \\

    \theutterance \stepcounter{utterance}  
    & & & \multicolumn{2}{p{0.3\linewidth}}{
        \cellcolor[rgb]{0.9,0.9,0.9}{
            \makecell[{{p{\linewidth}}}]{
                \texttt{\tiny{[GM$|$GM]}}
                \texttt{proposal} \\
            }
        }
    }
    & & \\ \\

    \theutterance \stepcounter{utterance}  
    & & & \multicolumn{2}{p{0.3\linewidth}}{
        \cellcolor[rgb]{0.9,0.9,0.9}{
            \makecell[{{p{\linewidth}}}]{
                \texttt{\tiny{[GM$|$GM]}}
                [1, 0, 0, 1]
            }
        }
    }
    & & \\ \\

    \theutterance \stepcounter{utterance}  
    & & & \multicolumn{2}{p{0.3\linewidth}}{
        \cellcolor[rgb]{0.9,0.9,0.9}{
            \makecell[{{p{\linewidth}}}]{
                \texttt{\tiny{[GM$|$GM]}}
                [[1, 0, 0, 1], [1, 0, 0, 1]]
            }
        }
    }
    & & \\ \\

\end{supertabular}
}

\end{document}
